% For compatibility with acroread version 
\pdfminorversion=4

\documentclass[mathserif, aspectratio=169]{beamer}
\mode<presentation>

%\usefonttheme[onlymath]{serif}
\usepackage[T1]{fontenc}
%\usepackage{tgheros}
%\usepackage{tgadventor}
%\usepackage{libertine}
%\usepackage{theinhart}
%\usepackage{palatino}
%\usepackage{newtxtext} %Times font only for text

% Use my theme
\usetheme{rognes}

% Used for image placement
\def\imagetop#1{\vtop{\null\hbox{#1}}}

% Misc packages
\usepackage{graphicx}
\usepackage[export]{adjustbox}
\usepackage{stmaryrd}
\usepackage{tikz}
\usepackage{booktabs}
\usepackage{multirow}
\usetikzlibrary{arrows,shapes,backgrounds,decorations,mindmap,patterns,snakes}
\usepackage{amsmath, amssymb, amsthm, amscd}
\usepackage[english]{babel}
\usepackage{algorithm}
\usepackage{pythonhighlight}
\usepackage{epstopdf}
\usepackage{multimedia}
\usepackage{soul}

\newcommand{\R}{\mathbb{R}}
\DeclareMathOperator{\Div}{\mathrm{div}}
\DeclareMathOperator{\Curl}{\mathrm{curl}}
\DeclareMathOperator{\Grad}{\mathrm{grad}}
\DeclareMathOperator{\D}{d}
\newcommand{\foralls}{\forall \;}
\newcommand{\iinner}[2]{\llangle #1, #2 \rrangle}
\newcommand{\inner}[2]{\langle #1,  #2\rangle}
\newcommand{\totald}{\mathrm{d}}
\newcommand{\refer}[1]{\begin{flushright}{\tiny \textcolor{darkgray}{[#1]}}\end{flushright}}
\newcommand{\referinline}[1]{{\tiny \textcolor{darkgray}{[#1]}}}
\newcommand{\ddt}[1]{{#1}_t}
\newcommand{\Mi}{M_i}
\newcommand{\Me}{M_e}
\newcommand{\Iion}{I_{\rm{ion}}}
\newcommand{\dx}{\, \mathrm{d}x}
\newcommand{\ds}{\, \mathrm{d} s}
\newcommand{\triang}{\mathcal{T}}
\DeclareMathOperator{\trace}{tr}
\newcommand{\duality}[2]{\langle #1, #2 \rangle}

\newcommand{\mysection}[1]{\begin{frame} \begin{center} \vspace{3em} \textbf{#1} \end{center} \end{frame}}
\newcommand{\videosection}[3]{\begin{frame} \begin{center} \vspace{3em} \href{#2}{\textcolor{rognesred}{\textbf{#1}}} \end{center} #3 \end{frame}}
\newcommand{\sectionwithrefer}[2]{\begin{frame} \begin{center} \vspace{3em} \textbf{#1} \end{center} \refer{#2} \end{frame}}

\newtheorem{prop}[theorem]{Proposition}


\tikzstyle{na} = [baseline=-.5ex]
\tikzstyle{every picture}+=[remember picture]
\everymath{\displaystyle}

\usepackage[small, figurename=,size=scriptsize,font={color=darkgray}]{caption}



%% \usepackage{listings}
%% \usepackage[most]{tcolorbox} 	%used for writing process notes, code boxes 
%% \tcbuselibrary{listings,skins}  % skins needed for code box styles
%% \usepackage{color}
%% \usepackage{amsfonts}
%% \usepackage{fancyvrb}		%provides tab and white space preservation for 
%% 				%the verbatim environment. Needed for Python
%% 				%code inclusion

\definecolor{programcode}{gray}{0.65}

\def\formtmpX#1#2{{\vskip3pt\noindent\fboxsep=0pt{\parbox{\textwidth}{\hbox to \textwidth{\hskip3pt\vbox{\raggedright\noindent\textbf{#2\vphantom{Qy}}}\hfill}}}\vskip3pt\par
\noindent\kern0pt}}

\newenvironment{programcode}[1]{\ignorespaces\def\stmtopen##1{##1}%
\formtmpX{programcode}{\centerline{\small{#1}}}}{\noindent\textcolor{programcode}{\rule{\columnwidth}{0pt}}\par\addvspace{\baselineskip}}%



\newcommand{\terminal}[1]{
  \vspace{-1em}
  \begin{programcode}{}%{Terminal window}%
    \colorbox{blue!10}{\parbox{0.98\textwidth}{\textcolor{black}{\texttt{#1}}}}
  \end{programcode}
  \vspace{-0.5em}
}

\newcommand{\emp}[1]{\texttt{#1}}

%-------------------------------------------------------------------------------
%-------------------------------------------------------------------------------

\DeclareMathOperator{\ssum}{\textstyle \sum}

\begin{document}

%-------------------------------------------------------------------------------
\cleanpage
\begin{frame}
  \begin{columns}
    \begin{column}{0.25\textwidth}
    \centering
    \vspace{1em}
    \includegraphics[width=\textwidth]{graphics/simula_logo_main_RGB.pdf} \\
    \vspace{1em}
    \includegraphics[width=\textwidth]{graphics/waterscales.pdf}  \\
    \vspace{1em}
    \includegraphics[height=0.3\textwidth]{graphics/flag_yellow_low.jpg} 
    \hspace{0.2em}
    \includegraphics[width=0.3\textwidth]{graphics/logo-erc.jpg} \\
    \vspace{1em}
    \includegraphics[width=0.5\textwidth]{graphics/waterscape_logo.png} \\
    \vspace{1em}
    \includegraphics[width=\textwidth]{graphics/rcn-logo.pdf} \\
    \end{column}
    \begin{column}{0.75\textwidth}
    {
    \centering

      \vspace{4em}

      {\bf Brain modelling: \\ from magnetic resonance images \\ to finite element simulation} \\

      \bigskip
      \bigskip
      

      Marie E. Rognes \\
      \bigskip
      \bigskip

      Oslo, Norway \\
      \medskip
      Jan 15, 22, 29 -- 2021 \\
      }
    \end{column}
  \end{columns}
\end{frame}

\cleanpage

\begin{frame}
  \frametitle{Outline of lectures}
  \begin{center}
  \includegraphics[height=2.3cm]{/home/meg/papers/current/hikeandsurf/Book/graphics/chp1/T1-image-rot-white}
  \includegraphics[height=2.3cm]{/home/meg/papers/current/hikeandsurf/Book/graphics/chp1/ernie-volume-64}
  \includegraphics[height=2.3cm]{/home/meg/papers/current/hikeandsurf/Book/graphics/chp1/soltn-t30-crop}
  \end{center}
  \begin{description}
    \item[Lecture 1-2] Quick start: From brain MRI to FEM
      \begin{itemize}
      \item
        Where can I find the slides (and book, data, and scripts)?
      \item
        Introduction to brain physiology and imaging
      \item
        From T1w MRI images to numerical simulation
      \end{itemize}
    \item[Lecture 3-4] Meshing hemispheres, brains, ventricles, parcellations
    \item[Lecture 5-6] Introducing anisotropy and heterogeneities (DTI)
  \end{description}
\end{frame}

\mysection{The brain as a poroelastic medium}

\input{/home/meg/presentations/slides/waterscape_poroelastic_brain.tex}
\input{/home/meg/presentations/slides/waterscape_volume_fraction_age.tex}
\input{/home/meg/presentations/slides/waterscape_volume_fraction_ecs.tex}

\mysection{What do we know of forces at play?}

\input{/home/meg/presentations/slides/waterscapes_respiratory_math_2019.tex}

\begin{frame}
\frametitle{Baledent}


\end{frame}

\mysection{Porous media brain modelling}

\begin{frame}
  \frametitle{Biot's equations describe displacement and flow in a poroelastic medium}
  \small
  \begin{columns}
  \begin{column}{0.6\textwidth}
    Find the displacement \only<1-3>{$u = u(x, t)$} and the pressure \only<1-3>{$p = p(x, t)$}
    over $\Omega \times [0, T]$ such that:
  \only<1-3> {
  \begin{subequations}
    \begin{align*}
    - \Div \left ( 2 \mu \varepsilon(u) + \lambda \Div u I - p I \right ) &= f, \\
    c_0 \dot{p} + \Div \dot{u} - \Div K \Grad p &= g
    \end{align*}
  \end{subequations}
  }

\vspace{-0.5em}
  \visible<2->{
  \textbf{Low-storage, incompressible regime}: $c_0 = 0$, $\lambda
    \rightarrow \infty$: $\Div u \rightarrow 0$, system decouples 
    \begin{align*}
      - \Div K \Grad p &= g &&\textrm{\alert{(Darcy)}} \\
      - \Div \left ( 2 \mu \varepsilon(u) + \lambda \Div u \right ) &= f - \Grad p && \textrm{\alert{(Elasticity)}} 
    \end{align*}
  }

  \vspace{-0.5em}
  \visible<3->{
    \textbf{Low-storage, impermeable regime}: $c_0 = 0$, $K \rightarrow 0$: 
      \begin{equation*}
        \begin{split}
        - \Div \left ( 2 \mu \varepsilon(u) + \lambda \Div u  - p I \right ) &= f, \\
        \Div u &= 0
        \end{split}
        \quad \textrm{\alert{(Stokes)}}
      \end{equation*}
      }
  \end{column}
  \begin{column}{0.4\textwidth}
    \centering
    \includegraphics[width=0.45\textwidth]{/home/meg/presentations/slides/graphics/mouse_u125.png} 
    \includegraphics[width=0.45\textwidth]{/home/meg/presentations/slides/graphics/mouse_p125.png} \\
    \vspace{-1em}
    \includegraphics[width=0.45\textwidth]{/home/meg/presentations/slides/graphics/mouse_v125.png} \\

    \includegraphics[height=0.2\textheight]{/home/meg/presentations/slides/graphics/biot.jpg}
    %\only<2>{\includegraphics[height=0.2\textheight]{/home/meg/presentations/slides/graphics/firth.jpg}}
    \visible<2->{\includegraphics[height=0.2\textheight]{/home/meg/presentations/slides/graphics/darcy.jpeg}}
    \hspace{1em}
    \visible<3->{\includegraphics[height=0.2\textheight]{/home/meg/presentations/slides/graphics/stokes.jpg}}
  \end{column}
  \end{columns}
  \vspace{-1em}
  \refer{Biot (1941), Murad, Thom\'ee and Loula (1992-1996), Phillips and Wheeler (2007-2008), and many others}
\end{frame}

\begin{frame}
  \frametitle{Stability of saddle point problems}
  {\small
    \vspace{-1em}
    \begin{columns}[T]
    \begin{column}{0.5\textwidth}
      \begin{block}{Abstract form}
      Find $u \in V$ and $p \in Q$ such that
      \begin{align*}
        a(u, v) + b(v, p) &= \inner{f}{v} \quad \foralls v \in V \\
        b(u, q) - c(p, q) &= \inner{g}{q} \quad \foralls q \in Q 
      \end{align*}

      \begin{equation*}
      \begin{pmatrix}
        A & B^T \\
        B & -C 
      \end{pmatrix}
      \begin{pmatrix}
        u \\
        p
      \end{pmatrix}
      =
      \begin{pmatrix}
        f \\
        g
      \end{pmatrix}
      \end{equation*}
      \end{block}
      \includegraphics[width=\textwidth]{/home/meg/presentations/slides/graphics/saddlepoint.png} 
      \refer{Ladyzhenskaya (1969),  Babu\v ska (1972), Brezzi (1974)}
    \end{column}
    \begin{column}{0.5\textwidth}
          \vspace{-1em}
      \textbf{Example: (Time-discrete) Biot}
        \begin{align*}
          A u &= - \Div (2 \mu \varepsilon(u) + \lambda \Div u I) \\
          B u &= - \Div u \\
          C p &= c_0 p - \Div K \Grad p
        \end{align*}
        \visible<2->{
        \centering
        \includegraphics[width=0.3\textwidth]{/home/meg/presentations/slides/graphics/cg.png}
        \includegraphics[width=0.3\textwidth]{/home/meg/presentations/slides/graphics/bdm.png}
        \begin{block}{Stability conditions}
        Well-posed problem if $a, b, c$ continuous, $c$ pos. def., and
        $\exists \, \alpha, \beta > 0$
        \begin{align*}
          a(u, u) \geq \alpha \| u \|_V^2 \quad \foralls u \in \mathrm{ker} B \\
          \inf_{q \in Q} \sup_{v \in V} |b(v, q)| \|v\|_{V}^{-1} \|q\|_{Q}^{-1} \geq \beta 
        \end{align*}
      \end{block}
        \vspace{-1em}
        }
    \end{column}
  \end{columns}
}
\end{frame}

\input{/home/meg/presentations/slides/waterscape_mpet_basics.tex}

\input{/home/meg/presentations/slides/waterscape_mpet_quasistatic.tex}


\mysection{Finite element methods for MPET}

\begin{frame}
  \frametitle{Preliminaries for variational formulation of the
    quasi-static MPET equations}

  \textbf{Boundary conditions:} Let $\Omega \subset \R^d$ with boundary
  $\partial \Omega$ and $T > 0$. \\ $u = 0$ on $\alert{\Gamma} \subset \partial
  \Omega$, $\sigma \cdot n \equiv (\sigma^{\ast}(u) - \ssum_{a}
  \alpha_a p_a I) \cdot n = 0$ on $\partial \Omega - \Gamma \not =
  \emptyset$ and \\ $p_a = 0$ on $\partial \Omega$ for $a = 1, \dots, A$.

  \medskip
  
  \textbf{Initial conditions}: $u(t = 0) = u_0$, $p_a(t = 0) = p_{a,
    0}$ for $a = 1, \dots, A$ ($c_a > 0$).

  \bigskip
  
  \textbf{Function spaces:} Define
  \begin{equation*}
    \alert{V} = H^1_{0, \Gamma}(\Omega; \R^d),
    \quad
    \alert{Q_0} = L^2(\Omega),
    \quad
    \alert{Q_a} = H^1_0(\Omega), \quad a = 1, \dots, A.
  \end{equation*}

  \bigskip

  \textbf{Weighted inner product and norm:} Let
  \begin{equation*}
    \inner{a}{b} \equiv \int_{\Omega} a \cdot b \dx,
    \quad
    \inner{a}{b}_{\beta} \equiv \int_{\Omega} \beta a \cdot b \dx,
    \quad
    \| a \|_{\beta}^2 \equiv \inner{a}{a}_{\beta}
    \quad
    \beta \in \R^{+} .
  \end{equation*}
\end{frame}


\input{/home/meg/presentations/slides/waterscape_soa_mpet.tex}

\input{/home/meg/presentations/slides/waterscape_total_pressure_short.tex}

\begin{frame}
  \frametitle{Displacements clearly differ between total pressure and "standard" mixed formulations ($\nu = 0.4999$)}
  \centering
  \begin{figure}
  \includegraphics[width=0.9\textwidth]{/home/meg/presentations/slides/graphics/waterscape_total_pressure.png}
  \caption{Comparison of displacements computed using the standard and
    total pressure formulation. Plots of displacement magnitude
    $|u(x_i, t)|$ versus time $t$, for a set of points $x_0, x_1,
    x_2$. The computed displacements clearly differ between the two
    solution methods}
  \end{figure}
\end{frame}

\mysection{}

\mysection{}

\mysection{}

\videosection{Video example 1}{}{}

\mysection{Igor, bring me more brains!}
 
\begin{frame}
\frametitle{Want to give it a try? Natural next steps:}
\begin{enumerate}
\item
\end{enumerate}
\end{frame}

%-------------------------------------------------------------------------------


\end{document}
